% Nguồn SGK Toán 9 Tập một — Luyện tập chung sau Bài 8, trang 52--53:
% https://cdn3.olm.vn/upload/thuvienso/4491668113-page-52-1773022937794.png
% https://cdn3.olm.vn/upload/thuvienso/4491668113-page-53-1773022938084.png

\section*{Luyện tập chung}

\begin{exercises}
    \textbf{Ví dụ 1}

    Không dùng MTCT, tính giá trị của biểu thức:
    \[
        A=\sqrt{(1-\sqrt{3})^2}-\sqrt{3}.
    \]

    \textbf{Giải}

    Ta có
    \[
        A=|1-\sqrt{3}|-\sqrt{3}=(\sqrt{3}-1)-\sqrt{3}=-1.
        \qquad \leftarrow \text{ hằng đẳng thức }\sqrt{A^2}=|A|.
    \]

    \textbf{Ví dụ 2}

    Tính:
    \begin{questions}
        \item $(2+\sqrt{5})(2-\sqrt{5})$;
        \item $(\sqrt{5}+\sqrt{2})(\sqrt{5}-\sqrt{2})$.
    \end{questions}

    \textbf{Giải}

    Sử dụng hằng đẳng thức hiệu của hai bình phương, ta được:
    \begin{questions}
        \item $(2+\sqrt{5})(2-\sqrt{5})=2^2-(\sqrt{5})^2=4-5=-1$.
        \item $(\sqrt{5}+\sqrt{2})(\sqrt{5}-\sqrt{2})=(\sqrt{5})^2-(\sqrt{2})^2=5-2=3$.
    \end{questions}

    \textbf{Chú ý.} Tổng quát, ta có
    \[
        (A+\sqrt{B})(A-\sqrt{B})=A^2-B
    \]
    với $A, B$ là hai biểu thức, $B\ge 0$.
    \[
        (\sqrt{A}+\sqrt{B})(\sqrt{A}-\sqrt{B})=A-B
    \]
    với $A, B$ là hai biểu thức không âm.

    Ta nói $A+\sqrt{B}$ và $A-\sqrt{B}$; $\sqrt{A}+\sqrt{B}$ và $\sqrt{A}-\sqrt{B}$ là những cặp biểu thức liên hợp.

    \textbf{Ví dụ 3}

    Cho biểu thức
    \[
        P=\dfrac{x\sqrt{x}+1}{x-\sqrt{x}+1}-1
    \]
    với $x\ge 0$.
    \begin{questions}
        \item Viết $x\sqrt{x}+1$ thành tổng hai lập phương rồi phân tích thành nhân tử để rút gọn biểu thức đã cho.
        \item Tính giá trị của $P$ tại $x=100$.
    \end{questions}

    \textbf{Giải}

    \begin{questions}
        \item Ta có
        \[
            x\sqrt{x}+1=(\sqrt{x})^3+1^3=(\sqrt{x}+1)\left[(\sqrt{x})^2-\sqrt{x}\cdot 1+1^2\right]
            =(\sqrt{x}+1)(x-\sqrt{x}+1).
        \]
        Từ đó
        \[
            P=\sqrt{x}+1-1=\sqrt{x}.
        \]
        \item Tại $x=100$ (thỏa mãn điều kiện) thì $P=\sqrt{100}=10$.
    \end{questions}

    \subsection{Bài tập}

    \begin{questions}[resume=SGK]
        \item Rút gọn các biểu thức sau:
        \begin{questions}
            \item $\sqrt{(\sqrt{3}-\sqrt{2})^2}+\sqrt{(1-\sqrt{2})^2}$;

            \answerline
            \answerline

            \item $\sqrt{(\sqrt{7}-3)^2}+\sqrt{(\sqrt{7}+3)^2}$.

            \answerline
            \answerline
        \end{questions}

        \item Thực hiện phép tính:
        \begin{questions}
            \item $\sqrt{3}\left(\sqrt{192}-\sqrt{75}\right)$;

            \answerline
            \answerline

            \item $\dfrac{-3\sqrt{18}+5\sqrt{50}-\sqrt{128}}{7\sqrt{2}}$.

            \answerline
            \answerline
        \end{questions}

        \item Chứng minh rằng:
        \begin{questions}
            \item $(1-\sqrt{2})^2=3-2\sqrt{2}$;

            \answerline
            \answerline

            \item $(\sqrt{3}+\sqrt{2})^2=5+2\sqrt{6}$.

            \answerline
            \answerline
        \end{questions}

        \item Cho căn thức $\sqrt{x^2-4x+4}$.
        \begin{questions}
            \item Hãy chứng tỏ rằng căn thức xác định với mọi giá trị của $x$.

            \answerline
            \answerline

            \item Rút gọn căn thức đã cho với $x\ge 2$.

            \answerline
            \answerline

            \item Chứng tỏ rằng với mọi $x\ge 2$, biểu thức $\sqrt{x-\sqrt{x^2-4x+4}}$ có giá trị không đổi.

            \answerline
            \answerline
        \end{questions}

        \item Trong Vật lí, tốc độ (m/s) của một vật đang bay được cho bởi công thức
        \[
            v=\sqrt{\dfrac{2E}{m}},
        \]
        trong đó $E$ là động năng của vật (tính bằng Joule, kí hiệu là J) và $m$ (kg) là khối lượng của vật (Theo sách Vật lí đại cương, NXB Giáo dục Việt Nam, 2016).

        Tính tốc độ bay của một vật khi biết vật đó có khối lượng $2{,}5$ kg và động năng $281{,}25$ J.

        \answerline
        \answerline
    \end{questions}
\end{exercises}
